\section{Conclusion}
% -----------------------------------------------------------------------------
A training observation cannot be reduced to a weight. Locally, it has a first-order component, the cotangent effort $\alpha_i$, and a positive second-order component, the observability tensor $\Gamma_i$, both relative to the collective geometry $G_S$. The triplet $(G_S,\Gamma_i,\alpha_i)$ is complete for the parametric action of the declared Gauss--Newton quadratic model, up to an additive constant, once the loss, output metric, collective geometry, and background metric have been fixed. An absolute utility may be completed by $c_i=\ell_i(\widehat\theta_S)$. On the principal stratum, the quotient has $2p$ global coordinates. More generally, an observation of rank $r<p$ belongs to a generic stratum of dimension $2r+1$, reduced to $2r$ when $b\in\Img B$. The compatibility lemma shows that the latter situation is natural for a quadratic loss with non-degenerate output curvature; the scalar case is therefore exactly the compatible $r=1$ stratum.

After whitening, this generator becomes a pair $(B,b)$ defined up to orthogonal action. The classification theorem shows that its orbit is fully determined by the spectrum of $B$ and the energy of $b$ in each eigenspace. Every reparameterization-invariant local quadratic quantity therefore factors through this signature. Leverage, information gain, Cook's distance, influence, and several marginal contributions appear as projections of this object.

The scalar limit is essential: when the observation is rank one, the invariant register reduces, for $q_i>0$, to the leverage--studentized-innovation-magnitude pair; the sign and invisible output components must be retained separately. Non-scalar structure becomes possible as soon as rank exceeds one; our simultaneous collision of five scalarizations is constructed in rank three and shows that these identical summaries are insufficient to recover the signature.

In the nonlinear setting, exact response depends on the Hessian whereas observability qualification relies on Gauss--Newton. Finite insertion follows a reweighting path and produces geometric feedback because the displacement changes the Jacobians and output metrics. Tracking $K_t^{(0)}$ and the representations confirms this mobility in the Friedman experiment without conflating raw parametric drift with functional change.

The same framework finally clarifies why label outliers, high-leverage points, invalid cells, and collection dependence call for different mechanisms. Displacement budgets, contraction budgets, structural validity, and statistical mass are distinct coordinates. The tabular prototypes illustrate this separation, but their nominal costs and calibration sensitivity confirm that they are not the central contribution of the paper; they belong to the tradition of GM-estimators.

The proposed contribution is therefore a local, conditional theory for qualifying observations before scalarization:
\[
\boxed{
\text{collective geometry}
+\text{point observability}
+\text{label effort}
=\text{contribution signature}.
}
\]
Its scientific scope does not lie in a universal victory over Ridge or Shapley, but in an intrinsic characterization of the information lost when this signature is replaced by a single score.

% -----------------------------------------------------------------------------
